%%%%%%%%%%  scan idx 306  |  folio 293  %%%%%%%%%%
elle définit donc un homomorphisme \(K^{\bullet}(X)_{\naif}\longrightarrow \varinjlim K^{\bullet}(X_i)_{\naif}\), dont on constate aussitôt, grâce à (EGA IV 8.5.2), qu'il est inverse de (3.2.1.1), ce qui prouve notre assertion.

\textbf{3.2.2.} On définit de même un homomorphisme canonique
\begin{equation}\tag{3.2.2.1}
 \varinjlim K^{\bullet}(X_i)\longrightarrow K^{\bullet}(X).
\end{equation}
Utilisant la technique de localisation de Deligne (SGA 4 XVII) qui implique que la catégorie des complexes parfaits sur \(X\) est équivalente à la catégorie \(\mathop{\underrightarrow{\mathrm{Lim}}}\) des catégories analogues \(\Parf(X_i)\), on trouve encore que (3.2.2.1) est un isomorphisme.

\textbf{3.2.3.} Supposons les \(X_i\) et \(X\) \emph{noethériens}, et les morphismes \(f_{ij}\) \emph{plats} (il en est alors de même des \(f_i\)). Alors les homomorphismes (2.12.1)
\(f_i^*:K_{\bullet}(X_i)\longrightarrow K_{\bullet}(X)\) définissent un isomorphisme
\begin{equation}\tag{3.2.3.1} \varinjlim K_{\bullet}(X_i)\xrightarrow{\sim}K_{\bullet}(X).
\end{equation}
La démonstration est analogue à celle de (3.2.1), en utilisant le fait (2.4) que \(K_{\bullet}(X)\simeq K_{\bullet}(X)_{\naif}\) se calcule en termes de faisceaux cohérents sur \(X\).

Comme les structures de modules des \(K_{\bullet}(X_i)\) sur les \(K^{\bullet}(X_i)_{\naif}\) sont compatibles avec les images inverses \(f_{ij}^*\) d'après (2.7 a)) et (2.12.2), on trouve, par passage à la limite, que \(\varinjlim K_{\bullet}(X_i)\) est muni d'une structure de module sur \(\varinjlim K^{\bullet}(X_i)_{\naif}\); celle-ci est compatible, via les isomorphismes (3.2.1.1) et (3.2.3.1) avec la structure de module de \(K_{\bullet}(X)\) sur \(K^{\bullet}(X)_{\naif}\), en d'autres termes on a un diagramme commutatif, où les flèches verticales sont définies par (3.2.1.1) et (3.2.3.1)~:
\[
\begin{tikzcd}[column sep=tiny]
\varinjlim K^{\bullet}(X_i)_{\naif}\otimes \varinjlim K_{\bullet}(X_i)
\arrow[r]\arrow[d,"\simeq"]
&\varinjlim K_{\bullet}(X_i)\arrow[d,"\simeq"]\\
K^{\bullet}(X)_{\naif}\otimes K_{\bullet}(X)\arrow[r]
&K_{\bullet}(X).
\end{tikzcd}
\]

